\chapter{Preface}

This book has a simple ambition: a single, coherent book of computer
science covering everything from the first year of school to the end of
the bachelor's degree, written in English for readers anywhere in the
world.

The style is concise and rigorous: courses built from definitions, examples,
theorems, methods and diagrams, with careful reasoning whenever it is
accessible at the given level, followed by carefully chosen exercises.
Algorithms come with their proofs of correctness and their costs. Results
stated without full proof are explicitly marked as admitted. Full solutions
to all exercises are collected at the end of the book. Every program printed
in these pages is taken from a file that has been run and tested, and every
output shown is the one the program actually printed.

The course comes in five books. The first covers the primary and middle
school years, the second the three years of high school, and the last
three the three years of university. The first years teach the machine
itself --- what a computer is made of, how to type with ten fingers, files,
memory and storage, the operating system, binary --- before the first
program is written. Programs are then written in Python throughout the
school years; the university books add C, C++, Rust and SQL. The two school
books never teach the same notion twice: each idea is taught once, in the
year where it first fits, and a later chapter that needs it opens with a
short reminder, \emph{You already know}, before going further. Where a
notion belongs to mathematics or physics, the companion books of the same
collection teach it, and this book recalls it before taking the view of the
computer scientist. In the school books each year is a part, split into
chapters; each university book is one year, split into parts and chapters.
The earlier the year, the younger the reader it is written for: more
illustrations, more detailed worked steps, and gentler exercises.

\section*{How to read this book}

Statements are numbered within each chapter and color-coded: definitions in
blue, theorems in red, propositions and their relatives in orange, and
methods --- short recipes for standard tasks --- in green. Programs are
printed with their line numbers and the name of their file; what a program
prints is shown below it on a grey background. Exercises are graded from
($\star$) for direct applications to ($\star\star\star$) for harder
problems. Framed boxes hold what is not part of the argument: a reminder of
what an earlier chapter taught, how professionals work in practice, a page
of history, facts that change from year to year (with the year they are
given for), and advice on safe use.

\section*{Safe use}

Everything this book asks of its reader happens on their own computer, with
free software, offline. Passwords, backups and updates get short boxes where
they belong. When a chapter explains how a program or a network can be
attacked, it does so on the book's own small programs and servers, running
on the reader's machine, so that the defence can be understood; attacking a
system that is not your own is never acceptable, and in most countries it is
a crime.

\section*{Contributing}

The source of this book lives in a public git repository:
\begin{center}
\url{https://github.com/bvirrion/one-computer-science-book}
\end{center}
Contributions --- new chapters, corrections, better explanations,
additional exercises --- are welcome; see \texttt{CONTRIBUTING.md} at the
root of the repository.
